Door Emmy’s defloratie
werd wat vast was gekraakt.
Al haar levensenergie
was na achttien jaar ontwaakt.

De deur van haar opening
zat niet langer meer op slot.
De slang van haar opwinding
lag niet langer in een knot.

Het kroop langs haar ruggengraat
naar steeds hogere sferen
om haar sluimerende staat
ook daar te activeren.

De hitte bevrijd door sex
ontdooide eensklaps het ijs.
Het triggerde de vortex,
de slinger van haar tijdreis.

De tijdreiziger keek op.
Een bewegend schilderij
herhaalde plaatjes non-stop.
Eronder wachtte een rij

van opgewonden mensen.
Ze vormden een lange haag.
Linten duiden de grenzen.
Daarbinnen kwamen gestaag

glimmende personen aan
uit de edele boten
van onyx of van mangaan,
die zonder zeil vlot vloten

en glansden na elk flitslicht.
De zwarte patrijspoorten
verborgen echter elk zicht
op de chique menssoorten,

die ze blijkbaar vervoerden.
Een rode loper lag uit.
Vele toeschouwers loerden
snel door een donkere ruit

als deze net niet dicht was.
Als de grootheid uitstapte
met een pompeuze grimas
en de menigte klapte,

verhitte adoratie
de gezwollen atmosfeer.
Ze leek sprekend op Emmy.
Zo droomde ze keer op keer

over deze knappe vrouw
van middelbare leeftijd
met dezelfde lichaamsbouw
waar men Emmy om benijdt.

“Graag een vochtvrije koffie!”,
zei een dikke kale man.
Op het naambord stond Eli.
Hij schonk water uit een kan,

omdat dat hem beter lag,
terwijl hij vergeefs wachtte
op een breed gedragen lach.
Waarop hij droog glimlachte:

“Expreszo dronk mijn maat die”
De rouwende menigte
kon niet wennen aan Eli
wat hun wrok verstevigde

over zijn onhandigheid.
Zijn afwijkende inborst,
die in een moeilijke tijd
wel vaker naast de pot morst.

Dan maar ernstig dacht Eli.
““Hoe lang moet jij nog” vroeg ie,
vaak uit levensallergie,
maar ook wel met ironie.”

Maar deze zware bromtoon
klonk nog meer uit de hoogte.
Zelfs in haar verdriet beeldschoon
schudde haar hergeboorte

- zoals Emmy de vrouw zag -
uit gêne van links naar rechts.
Ze ervoer zijn wangedrag
voor het kritisch publiek slechts

als plaatsvervangende schroom.
Hij was immers haar pronkstuk
en vroeger haar meisjesdroom.
Een toegangspoort naar geluk.

Hij was ooit slank, knap en wijs,
maar werd vet, oud en gênant.
Hij kocht haar een paradijs
op een fraai Waddeneiland.
Een chique kerk naar haar smaak
verbouwd met een ruim budget.
Al haar pijlen schoten raak,
maar haar geest verklootte het.

De eeuwig kritische stem,
die haar in de hel plaatste.
De schuld legde ze bij hem.
Haar onvrede weerkaatste.

Ergo kropen haar schouders
steeds vaker bij hem omhoog.
Net als ooit bij haar ouders
was het angst dat haar leeg zoog.

Het verlies van zekerheid
in haar partners nabijheid,
leidde tot nervositeit,
conflict en verlegenheid.

“Je bent toch psychiater?!
Had dan onze vriend gered!
In plaats van hier je flater
te dichten in je gebed.

In plaats van eerbied en spijt
voer je weer het hoogste woord.
Zijn spiritualiteit
leidde juist tot zijn zelfmoord.

Zijn geloof in verlossing
heeft het motief geschapen.
De leegte als oplossing
leverde het moordwapen.

Ons verdriet door het verlies
bagatelliseer je bot
met jouw pedant rouwadvies.
Het universum als God

zou hem laten herleven
volgens jouw overtuiging
om meer inzicht te geven
voor een hernieuwde poging.

Dus hij is volgens jouw leer
even op een tussenreis
en bezoekt ons keer op keer,
maar dan steeds meer levenswijs

tot de keten van zijn levens
zijn ziel heeft getransformeerd.
Daarmee beweer je tevens
dat hij te zwak heeft geleerd

in zijn huidige bestaan.
Ik vind dat aanmatigend.
Hij is er vol voor gegaan.
Het bleek overweldigend.

Echter waar echt de schoen wringt
is dat je ons bij zijn uitvaart
die filosofie opdringt.
Vond hij die echt zoveel waard?”

Emmy’s dubbelgangster gaf
haar man geen enkele steun,
maar liet vergeldend en laf
hem bungelen als miskleun.

Hij verdedigde zich niet
ook al kon hij ruim putten
uit zijn eigen vakgebied.
Hij wilde niet opjutten,

maar zijn vakkennis delen.
Bijna-dood-ervaringen
waren de bestanddelen
onder zijn verklaringen

zonder gelijk te claimen.
Hij hielp alleen cliënten
met psychische problemen
door zulke accidenten.

Fan zegt: “Maak misbruik van me!”
Celeb zegt: “Wie ben je dan?”
Fan zegt: “Jouw slavinnetje!”
Celeb zegt: “Ik geloof er niets van.

Dan had je geen kleding aan.”
Op een geletterd klavier
ziet Emmy vingers stilstaan.
Stralend glas was het papier

waar de toetsen hun tonen
als letters lieten schijnen.
Geen lood of inktpatronen.
Verschijnen en verdwijnen

door een zachte vingertik.
De handen vol met sproetjes
aarzelden vanwege schrik.
Ze wiebelde haar voetjes

onder een eiken mensa.
Fan zegt: “Het spijt me, patron!
U wilt dat ik verderga
dan ik eerder voor u kon?”

Een duim kwam als tekening
op het levend schilderij.
Weer was er de aarzeling.
Toch stapte de vrouw opzij

schuivend van haar rieten stoel
en bewoog zich van het doek,
weloverwogen en zwoel
in hemd zonder onderbroek.

Emmy herkende de vrouw.
Ze bleek haar zielsverwante.
Dit keer was ze braaf en trouw
en niet een bijdehante.

Een rood lichtje in de lijst
leek ineens te knipperen.
“Ik wil dat je het ophijst!”
Ruimte om te schipperen

leek de tekst niet te geven,
dus ging haar singlet omhoog.
Al deed ze dit maar even
toch bleek dat ze hem bedroog

met wie ze was verbonden.
Een fontein zonnestralen
kleurde haar wulpse zonden.
Bleef het eerst bij verhalen,

kwamen er plots beelden bij.
De mystieke schaduwen
op linker- en rechterdij
kwamen over als kluwen

van een lubrieke plaaggeest,
die in haar oranje nest
uitgebreid hadden gefeest.
De fan slaagde voor zijn test.

Emmy vloog ondertussen
van de ook roodgetinte,
wiens ogen scherp focussen
op hetgeen celeb printte,

naar de glas-in-lood ramen,
die de regenboog straalden
op de vrouw haar examen.
Emmy’s kraalogen dwaalden

via het stoffig zonlicht
naar een lijst boven de schouw
van de authentieke haard,
die Emmy’s kopie graag wou.

De vrouw bleek erin gelegd,
maar nog jong en verwilderd.
Het portret leek levensecht.
Zo scherp was het geschilderd.

Het beeld begon te praten
en bewoog tot Emmy’s schrik.
Ze werd achtergelaten
onder protest en gesnik

door haar ouders op een plein,
met een zandbak en klimrek.
Er bleek geen erger venijn
dan die van een lachebek

happend op kwade tongen.
Haar beste vriendin deed mee.
“Menseshaar!”, riep een jongen.
Het was zijn nieuwe trofee.

Het beschermde de schavuit
tegen de autoriteit.
“Latijn voor ongesteldheid”,
legde hij fluisterend uit

aan de vragende pestclub.
“maar het klinkt als mensenhaar!”
Een ander bedacht “Ketchup”,
dat schalde wat minder raar.

“Wie helpt haar?”, dacht Emmy.
“Zijn rode haren kinky?”,
vroeg huilend haar jeugdversie
ineens staand voor haar mammie.

“Daar kun je toch niets aan doen.
Zo ben je door God verdoemd.”
Waarna de tactloze oen
slechts haar eigen jeugd benoemd,

waarin zij steeds werd gejend.
Hoe weer haar moe zich schaamde
en dat aan haar had bekend.
Daarmee de vlek beaamde.

Slim het slachtoffer gespeeld
scoorde de emotiedief.
Emmy’s jonge evenbeeld
troostte maar haar moederlief.
